\documentclass{article}

\usepackage[margin=1in]{geometry}

\usepackage{parskip}

\begin{document}

% Page 1

Combining Valence-to-Core X‑ray Emission and Cu K‑edge X‑ray
Absorption Spectroscopies to Experimentally Assess Oxidation State
in Organometallic Cu(I)/(II)/(III) Complexes
Blaise L. Geoghegan, Yang Liu, Sergey Peredkov, Sebastian Dechert, Franc Meyer, Serena DeBeer,
and George E. Cutsail III*
Cite This: J. Am. Chem. Soc. 2022, 144, 2520−2534
Read Online
ACCESS
Metrics \& More
Article Recommendations
*
sı
Supporting Information
ABSTRACT: A series of organometallic copper complexes in formal oxidation
states ranging from +1 to +3 have been characterized by a combination of Cu K-
edge X-ray absorption (XAS) and Cu Kβ valence-to-core X-ray emission
spectroscopies (VtC XES). Each formal oxidation state exhibits distinctly
different XAS and VtC XES transition energies due to the differences in the Cu
Zeff, concomitant with changes in physical oxidation state from +1 to +2 to +3.
Herein, we demonstrate the sensitivity of XAS and VtC XES to the physical
oxidation states of a series of N-heterocyclic carbene (NHC) ligated
organocopper complexes. We then extend these methods to the study of the
[Cu(CF3)4]−ion. Complemented by computational methods, the observed
spectral transitions are correlated with the electronic structure of the complexes
and the Cu Zeff. These calculations demonstrate that a contraction of the Cu 1s orbitals to deeper binding energy upon oxidation of
the Cu center manifests spectroscopically as a stepped increase in the energy of both XAS and Kβ2,5 emission features with increasing
formal oxidation state within the [Cun+(NHC2)]n+ series. The newly synthesized Cu(III) cation [CuIII(NHC4)]3+ exhibits
spectroscopic features and an electronic structure remarkably similar to [Cu(CF3)4]−, supporting a physical oxidation state
assignment of low-spin d8 Cu(III) for [Cu(CF3)4]−. Combining XAS and VtC XES further demonstrates the necessity of combining
multiple spectroscopies when investigating the electronic structures of highly covalent copper complexes, providing a template for
future investigations into both synthetic and biological metal centers.
■INTRODUCTION
Copper’s ability to facilitate a wide-range of catalytic processes
has made it a highly popular metal for homogeneous catalysis.
Copper has been successfully employed as a catalyst in C−C
bond formation, C−H bond activation and is especially
prevalent in biological oxidation and oxygenation processes.1
Many of the proposed pathways for homogeneous copper-
mediated reactions invoke redox couples including not only
CuI and CuII but also the CuIII ion;2−6 therefore, under-
standing the spectroscopic markers that distinguish these three
physical oxidation states of Cu from one another is important
in order to determine reaction mechanisms in both synthetic
and biological settings. Within the literature, the assignment of
a Cu(III) oxidation state to catalytic intermediates7 has been
met with a degree of controversy due to the decreasing
calculated metal 3d admixture in the frontier molecular orbitals
(MOs) of high-valent late transition metals.8−10 For instance,
the physical oxidation state and electronic structure of
[Cu(CF3)4]−anion has been debated for decades.10−14 The
diamagnetic and colorless complex has been structurally and
spectroscopically characterized as a closed-shell S = 0 species
with distorted D2d symmetry. One may describe this complex
as it was originally done: a formal CuIII center ligated by four
monodentate, monoanionic CF3
−ligands.15 Alternatively, a d10
(CuI) electronic configuration can be assigned, where the
ligands assume either a heterogeneous 3(CF3
−) + CF3
+ or
homogeneous 4(CF3
0.5−) oxidation state.11,16 To assign a
formal Cu(I) oxidation state, a select definition of an “inverted
ligand field” is used, where the MOs with >50\% Cu 3d
admixture sink to energies below those with mostly ligand
admixture. The resultant lowest unoccupied molecular orbital
(LUMO) in this picture is therefore predominately of ligand
parentage and a supposed 3d10 electronic configuration for
[Cu(CF3)4]−is thus invoked.11 The electronic structure of
copper centers like the [Cu(CF3)4]−anion have led to
unique14 and differing17 spectroscopic interpretations.
As an atom-specific probe by the means of core electron
excitation, X-ray absorption spectroscopy (XAS) is frequently
Received:
September 8, 2021
Published: January 20, 2022
Article
pubs.acs.org/JACS
© 2022 The Authors. Published by
American Chemical Society
2520
https://doi.org/10.1021/jacs.1c09505
J. Am. Chem. Soc. 2022, 144, 2520−2534
Downloaded via FREIE UNIV BERLIN on November 28, 2025 at 14:04:26 (UTC).
See https://pubs.acs.org/sharingguidelines for options on how to legitimately share published articles.

\newpage

% Page 2

employed to determine oxidation states.18−20 The edge energy
shifts in both K (1s) and L (2s,p) edge XAS or X-ray
photoemission spectroscopy (XPS) often correlate well to the
core orbital ionization energy.9,21−30 For copper, XAS has
proven itself a reliable probe of copper’s oxidation
state19,20,31,32 and has been used to show large shifts in
spectra for Cu(I), Cu(II), and Cu(III) species.18,20 For
formally assigned Cu(III) species, consistent shifts to higher
energies have been observed in both K and L-edge XAS,
indicating an increased effective nuclear charge (Zeff) at the Cu
upon oxidation. However, decreased L-edge intensity of formal
Cu(III) centers has been reinterpreted to support a Cu(I)
physical oxidation state assignment where an inverted field is
present.10 This is in contrast to the previously established
interpretation of increased metal−ligand covalency for the high
oxidation-state, yielding decreased copper 3d character in the
frontier orbitals.28
We previously reported the high-resolution Cu K-edge XAS
of a series of copper +1, +2, and +3 complexes ligated by a
neutral, tetradentate 16-atom macrocyclic ligand with trans N-
heterocyclic carbene (NHC) donors and trans pyridine donors
(Scheme 1).32 Here, the flexibility of the pyridine donors of
the macrocycle allows the Cu(I) species to adopt an
approximately linear two-coordinate environment, whereas in
the Cu(II) and Cu(III) centers both the pyridyl and NHC
groups coordinate to the Cu center with a roughly square-
planar CNCN first coordination sphere, accompanied by a
distant (weakly bound) apical solvent (MeCN) ligation. The
strong NHC σ donors of the ligand and their high covalency
with the copper center are speculated to aid in the stabilization
of the high-valent Cu(III) ion. Supported by X-ray
crystallography, UV−vis spectroscopy, electrochemical charac-
terization, and theoretical approaches, the physical oxidation
states of copper in this redox series were clearly assigned via
shifts in the energy of the pre-edges and edge features of their
Cu K-edge XAS.32
The paramagnetic nature of Cu(II) complexes allows for
their extensive characterization through EPR and/or nuclear
resonance spectroscopy, which provides a rich picture of both
the electronic and geometric structures. However, for the EPR
silent Cu(I) and Cu(III) species, complementary techniques to
XAS that offer additional insight into the electronic structure of
the Cu centers are limited. X-ray emission spectroscopy (XES)
results from photon emission due to the relaxation of electrons
from orbitals with principle quantum numbers n > 1 into the 1s
core hole produced by a high-energy incident X-ray beam
(Figure 1).33 The Cu Kβ XES spectrum is comprised of two
main regions of interest. At lower energy, the mainline region
(∼8905 eV) originates from the refilling of the core 1s hole by
Cu 3p electrons.25,34 It is well demonstrated for other 3d (and
4d) transition metals that the XES mainlines are excellent
reporters of 3d (or 4d) electronic structure due to the role of
p-d electron spin exchange interactions on the spectrum.34−36
The mainlines of both small molecule and protein-bound
Cu(I) centers have been more recently reported.37 Resonant
inelastic X-ray scattering (RIXS) measurements have enabled
more detailed analysis of less featured Cu mainlines, leading to
distinct assignments of charge-transfer and multielectron
transitions achieved at higher resonant excitation energies,
corresponding to previously proposed metal−ligand charge
transfer (MLCT) features of Cu(I) centers.38 Classically,
nonresonant Kβ mainline spectra have been used for the
fingerprinting of transition metal spin states;39,40 however,
covalency has been shown to complicate this simple
picture.35,41 We note, however, that Cu Kβ mainlines have
yet to be thoroughly explored as reporters of oxidation and/or
spin state, although the copper mainline is reasoned to follow
the same qualitative trends as other dn > 5 transition metals.
At higher energy, but with lowest transition probability and
oscillator strength, the valence-to-core (VtC) region (∼8975
eV) is subdivided into the Kβ″ and Kβ2,5 features, originating
from the refilling of the core 1s hole by MOs of predominantly
ligand ns and np character, respectively.33,42−44 Hence, VtC
XES directly probes the energy of the valence MOs with
respect to the metal 1s orbital and has been shown to be highly
sensitive to ligand type,38,45−47 bond activation,45,48,49
protonation state,50,51 and ligand orientation.37,38,49 Despite
the fact that VtC spectroscopy has been shown to be sensitive
to the oxidation state in other first-row transition metals such
as Mn39 and Fe,52 it has not yet been employed to aid in the
determination of oxidation state in Cu centers.
For the electronic structure of organocopper complexes, we
hypothesized that VtC XES will offer new spectroscopic insight
for understanding of copper’s electronic structure. It has been
demonstrated that strong metal−ligand covalency can
complicate oxidation state assignments based on Cu K-edge
XAS alone.18 Therefore, employing multiple core X-ray
spectroscopies to interrogate the electronic structure of such
high-valent Cu systems may provide observable trends that are
related to physical oxidation states. These spectroscopic
observables form the foundations for such analysis of highly
covalent systems, for which computational methods may give
an ambiguous picture.
Scheme 1. Organometallic Copper Complexes Investigated
in the Current Work
Journal of the American Chemical Society
pubs.acs.org/JACS
Article
https://doi.org/10.1021/jacs.1c09505
J. Am. Chem. Soc. 2022, 144, 2520−2534
2521

\newpage

% Page 3

Given the highly covalent nature of both NHC and CF3
−
ligands, we present here the series of NHC/pyridine
macrocycle complexes as a calibrant for the use of VtC XES
as an additional experimental technique in conjunction with
Cu K-edge XAS in the determination of physical oxidation
state and for understanding the nature and origins of the
observed transitions. Additionally, for a systematic extension of
the NHC-bearing macrocyclic systems toward the homoleptic
C4 first coordination sphere, more analogous to [Cu(CF3)4]−,
a
new
tetra-NHC
Cu(III)
macrocyclic
complex
[CuIII(NHC4)]3+ has been prepared and included in this
study. [CuIII(NHC4)]3+ is similar to the tetra-NHC Cu(III)
complex recently reported by the Kühn group,53 but has
backside methyl substituents at the four imidazol-2-yliden
groups (Scheme 1). The relationships and trends extracted
from the investigation of the four NHC-bearing complexes are
then applied to both the VtC XES and Cu K-edge XAS of the
[Cu(CF3)4]−complex in order to determine its physical
oxidation state and electronic structure. The combination of
the complementary XAS and VtC XES techniques is shown to
be crucial for assessing the complicated electronic structures of
highly covalent Cu complexes, overcoming the limitations of
analyzing results from either technique in isolation. Through-
out this report, we will emphasize the importance of
correlating observables across multiple measurement techni-
ques and how this forms a firm foundation for interpretation of
complicated electronic structures, especially in highly covalent
systems.
■RESULTS
The complexes of general formula [Cun+(NHC2)]n+ form a Cu
redox series of a macrocyclic ligand bearing trans carbene and
trans pyridyl donor groups with a central Cu(I)/(II)/(III)
ion.32 It was previously shown that the redox state of the ligand
Figure 1. Simplified MO diagram schematically depicting the origin of the transitions in the Cu Kβ XES spectrum (left) and the experimental Cu
Kβ XES spectrum of [CuIII(NHC4)]3+ with VtC region enlarged approximately 100-fold (right).
Scheme 2. Synthesis of the New Tetra-NHC Ligated Cu(III) Complex (left) and the Molecular Structure of Its Cation
[CuIII(NHC4)]3+ in Solid State (Right; Side View and Top View)
Journal of the American Chemical Society
pubs.acs.org/JACS
Article
https://doi.org/10.1021/jacs.1c09505
J. Am. Chem. Soc. 2022, 144, 2520−2534
2522

\newpage

% Page 4

remains fixed throughout the series, but in [CuII(NHC2)]2+
and [CuIII(NHC2)]3+ an acetonitrile molecule is weakly bound
in an apical site.32 The new complex [CuIII(NHC4)]3+ is
closely related to [CuIII(NHC2)]3+, but the two pyridyl donors
are substituted for NHC groups and the backbones of the
imidazole-2-ylidenes are methylated. [CuIII(NHC4)]3+ was
synthesized according to the method that was originally
described by Ghavami et al.53 for the nonmethylated parent
analogue, with slight modifications (Scheme 2, left; see
Supporting Information for details of the synthetic protocols).
The PF6
−salt of [CuIII(NHC4)]3+ was fully characterized by
various spectroscopic methods, CHN analysis as well as by
single crystal X-ray diffraction (see SI for detailed information
and spectroscopic data; crystallographic parameters and
refinement details are given in Table S2, and selected bond
lengths and angles for the [CuIII(NHC4)]3+ cation are given in
Table S3). The molecular structure of [CuIII(NHC4)]3+ in the
solid state (Scheme 2 and Figure S16) reveals no binding of
acetonitrile in the apical position, although two acetonitrile
solvent molecules are present in the asymmetric unit of the
crystal lattice. The primary C4 coordination sphere in
[CuIII(NHC4)]3+ is almost perfectly square planar (sum of
angles around Cu is 360.0°), and the Cu−CNHC bond
distances are in the range 1.877(3)−1.889(3) Å (average
1.88 Å) and thus very similar to those of [CuIII(NHC2)]3+
(average 1.88 Å). The overall macrocycle conformation is
nonplanar and saddle-shaped, but variable-temperature NMR
(VT-NMR) spectra show facile conformational inversion of
the macrocyclic ligand in MeCN-d3 solution (with activation
parameters ΔH‡ = 17.8 ± 0.2 kcal mol−1, ΔS‡ = 19.0 ± 0.8 cal
mol−1 K−1 derived from an Eyring plot; see SI for details) with
apparent D4h symmetry on the NMR time scale at room
temperature. NMR data (Figures S2−S11) and SQUID
magnetometry (Figure S15) confirm that [CuIII(NHC4)]-
(PF6)3 is diamagnetic over the entire temperature range
studied.
Cu K-edge XAS of NHC Complexes. The Cu K-edge
XAS is presented initially as a metric for both physical
oxidation state of the copper center, as well as the metal−
ligand covalency. The Cu K-edge Kβ high-energy resolution
fluorescence detected (HERFD) XAS spectra of the
[Cun+(NHC2)]n+ series has been discussed in detail
previously.32 In parallel to those measurements, the more
typical transmission mode XAS data were also collected and
are now displayed in Figure 2a and Figure S17. Transmission
mode Cu K-edge XAS was obtained for [CuIII(NHC4)]3+ and
[Cu(CF3)4]−at the SAMBA beamline at the SOLEIL
synchrotron (see SI for sample and collection details), and
the results are compared to the [Cun+(NHC2)]nn series (Figure
2a). The spectrum of the [Cu(CF3)4]−species is a very good
match to that previously reported.14
The [CuI(NHC2)]+ complex exhibits an absorption profile
unlike the Cu(II) and Cu(III) species due to its d10 electronic
configuration, and thus lack of a classic 1s →3d pre-edge
transition. The lowest energy feature in the XAS of
[CuI(NHC2)]+ is an electric dipole allowed 1s →4pz
transition at 8982.0 eV and is followed by an even more
intense feature at 8985.1 eV, which is likely attributed to a 1s
→4py excitation based on molecular symmetry (Scheme 1).32
To higher energy is a poorly resolved feature at 8989.5 eV, and
the white line maximum is observed at ∼8994 eV.
[CuII(NHC2)]2+ and [CuIII(NHC2)]3+ both exhibit formally
electric dipole forbidden 1s →3d pre-edge transitions at
∼8979.5 and 8981.3 eV, respectively. While [CuIII(NHC4)]3+
lacks an isolated pre-edge feature of similar energy to its
[CuIII(NHC2)]3+ analogue, it does exhibit a more intense
shoulder at ∼8982.3 eV, approximately 1 eV higher than
[CuIII(NHC2)]3+, that is most clearly observed in the first
derivative spectrum shown in Figure 2b, which can be
putatively assigned as a 1s to 3d pre-edge transition. This is
slightly higher than the typical 8981 ± 0.5 eV range observed
for other Cu(III) systems20 and suggests that the exchange of
pyridyl donors for strong σ donating carbenes in the
Figure 2. (a) Cu K-edge XAS (transmission mode) spectra for all
compounds (solid lines). (b) Expansion of pre-edge region (left) and
first derivatives of experimental and TDDFT calculated (B3LYP/
ZORA-def2-TZVP/ZORA) pre-edge region for the [CuII(NHC2)]2+,
[CuIII(NHC2)]3+, [CuIII(NHC4)]3+ species. Vertical dashed lines are
included as visual guides. (c) Transition difference density plots for
the pre-edge transition in [CuIII(NHC2)]3+ and [CuIII(NHC4)]3+
plotted at an isosurface value of 0.003 au. Legend: C, gray; Cu,
orange; N, blue. Hydrogen atoms omitted for clarity.
Journal of the American Chemical Society
pubs.acs.org/JACS
Article
https://doi.org/10.1021/jacs.1c09505
J. Am. Chem. Soc. 2022, 144, 2520−2534
2523

\newpage

% Page 5

macrocyclic tetra-NHC ligand significantly influences the 1s-
LUMO gap (antibonding interaction with the Cu 3dx
2−y
2) in
the
resulting
Cu(III)
complex
(vide
inf ra).
The
[CuIII(NHC4)]3+ spectrum is highly featured and exhibits
both similarities and differences to that of the bis-NHC
analogue [CuIII(NHC2)]3+. [CuIII(NHC4)]3+ exhibits three
main, intense features in the rising edge at approximately
8985.0, 8989.2, and 8993.0 eV, with a shoulder at ∼8997.0 eV
before the white line is observed at ∼9005 eV. Both
[CuIII(NHC2)]3+ and [CuIII(NHC4)]3+ exhibit similar edge
positions,
which
are
shifted
∼5
eV
higher
than
[CuII(NHC2)]2+, supporting increased oxidation state assign-
ments as the effective 1s binding energy increases for the
[CuIII(NHC2)]3+ and [CuIII(NHC4)]3+ relative to the Cu(II)
complex. Relative to that of the Cu(I) species, the Cu(II) and
Cu(III) NHC complexes exhibit white lines of both higher
energy and higher intensity, in line with the increased
oxidation state: [CuII(NHC2)]2+ 8996.5 eV; [CuIII(NHC2)]3+
9000.0 eV; [CuIII(NHC4)]3+ 9005.0 eV. However, differences
between the ligands (i.e., pyridine vs carbene) may also play a
role here, as [CuIII(NHC2)]3+ exhibits a slightly sharper lower
energy white line than its [CuIII(NHC4)]3+ analogue. Apart
from the higher energy pre-edge transition in [CuIII(NHC4)]3+
compared to [CuIII(NHC2)]3+, the main difference between
the two spectra is that the [CuIII(NHC4)]3+ exhibits an intense
and well-resolved feature at ∼8985.0 eV, isoenergetic with the
most intense rising edge feature in the Cu(I) species
[CuI(NHC2)]+, whereas [CuIII(NHC2)]3+ does not. However,
the intensity of the 8985 eV feature of the [CuIII(NHC4)]3+
spectrum is significantly less than the Cu(I) feature and is
comparable to the previously assigned 1s →4p features of
[CuIII(NHC2)]3+.32 The energy of both the pre-edge and white
line features of the [CuIII(NHC4)]3+ complex are distinctly
different from that of the [CuI/II(NHC2)]+/2+ species, but with
sufficient similarities to [CuIII(NHC2)]3+ in order to make a
Cu(III) assignment.
As previously demonstrated for [CuII(NHC2)]2+ and
[CuIII(NHC2)]3+, TDDFT calculations of the Cu K-edge
spectrum exhibit a single pre-edge transition into an MO that
is clearly of copper 3dx
2−y
2 parentage and their energy
differences are consistent with the deeper binding energy of
the Cu 1s electrons in the higher valent systems.32 Extending
these calculations to [CuIII(NHC4)]3+ predicts that the pre-
edge transition in tetra-NHC complex [CuIII(NHC4)]3+ is
shifted into the rising edge by ∼0.9 eV compared to the
equivalent transition in [CuIII(NHC2)]3+, lending support to
the notion that modulation of the LUMO energy is achieved
through increased covalency. In fact, [CuIII(NHC2)]3+ and
[CuIII(NHC4)]3+ exhibit nearly identical transition state
difference densities, as shown in Figure 2c. It is clear from
the inspection of the XAS alongside the TDDFT that the
classic 1s →3dx
2−y
2 transition for [CuIII(NHC4)]3+ is quite
high in energy, even approaching the edge structure. The
observed trend of increased pre-edge energies with increased
metal−ligand covalency is in line with other organocopper
complexes, and the substitution of C-donor atoms into the first
coordination sphere has been demonstrated to aid in the
stabilization of the Cu(III) ion.54 TDDFT calculations also
predict with good accuracy the energy separation of the pre-
edge feature and the Cu 1s →4pz and Cu 1s →4px,y
transitions in [CuIII(NHC4)]3+ as shown in Figure S19.
Cu K-edge XAS of [Cu(CF3)4]−. Overall, the absorption
profile of the [Cu(CF3)4]−complex has some similarities with
those of the Cu(III) macrocyclic NHC complexes, especially
the [CuIII(NHC4)]3+ complex, despite fundamentally different
ligand environments. The pre-edge transition in [Cu(CF3)4]−
occurs at ∼8982 eV (Figure 2 and Figure S19), similar to that
observed in [CuIII(NHC4)]3+ and slightly higher in energy
than in [CuIII(NHC2)]3+. This suggests a correlation between
higher metal−ligand covalency and increased energy of the
pre-edge transition, implying that the four CF3
−ligands form a
more covalent interaction with the Cu center than only two
NHCs and two pyridines, but perhaps not as covalent as the
four NHCs of the [CuIII(NHC4)]3+ complex.
To higher energy, the 8985.6 eV feature exhibited by
[Cu(CF3)4]−has been ascribed to Cu 1s →4pz excitations,14
and is more intense than the ca. 8985 eV transition of the
[CuIII(NHC4)]3+ species, which has a similar 1s →4pz
transition (Figures S19−20 and Table S4). The less
pronounced features in the rising edge of [Cu(CF3)4]−at
∼8989.5 and 8994.0 eV generally have nearly isoenergetic
counterparts
in
both
the
[Cu III(NHC2)]3 +
and
[CuIII(NHC4)]3+ species ∼0.3 eV lower in energy. Addition-
ally, the white line in the [Cu(CF3)4]−complex is intermediate
to that of [CuIII(NHC2)]3+ and [CuIII(NHC4)]3+, mirroring
the trend in the pre-edge energies and thus the potential
covalency trend. Despite the similarities in the spectra of the
[Cu(CF3)4]−and [CuIII(NHC4)]3+ complexes, there are key
differences when compared to the [CuIII(NHC2)]3+ complex:
chiefly, the much more intense 1s →4pz transition at ∼8985
eV and the much lower intensity feature at 8989.5 eV. This is
unsurprising due to the differences in ligand types and
molecular geometries (Scheme 1). Noteworthily, the pre-
edge feature and 1s →4pz feature of [Cu(CF3)4]−are
isoenergetic with the lowest energy and the most intense
absorption features of the Cu(I) species [CuI(NHC2)]+,
respectively. However, the low intensity of the pre-edge
feature and higher energy rising edge and white line of the
[Cu(CF3)4]−species relative to that of the Cu(I) species
suggest significant d-hole character as well as a significantly
increased Cu 1s ionization energy, further supporting a Cu(III)
assignment based on the experimental data.18
TDDFT calculations reproduce the Cu K-edge XAS for
[Cu(CF3)4]−, showing that the pre-edge transition is due to an
excitation of a core 1s electron into what is clearly a dx
2−y
2
localized Cu orbital with σ bonding to the CF3
−ligands
(Figure S19). The calculated pre-edge transition difference
density plots for [Cu(CF3)4]−and [CuIII(NHC4)]3+ exhibit
remarkably similar characteristics, suggesting a similar LUMO
electronic structure of the two complexes. Additionally, the
LUMO in [Cu(CF3)4]−transforms as b2 in approximate D2d
symmetry, meaning that mixing of the Cu 4pz orbital into the
dx
2−y
2 is allowed by symmetry, thus increasing the intensity of
the pre-edge feature through reinstating partial electric dipole
allowed character to the transition.
The clear assignment of the Cu(I), (II) and (III) oxidation
states in the [Cun+(NHC2)]n+ series via their Cu K-edge XAS
can be used here as a guide for evaluation of the oxidation state
of the Cu ions in both the [CuIII(NHC4)]3+ and [Cu(CF3)4)]−
complexes. In combination with the crystallographic character-
ization, the Cu K-edge XAS strongly supports the high valent
+3 oxidation state of Cu in the [CuIII(NHC4)]3+ species.
Additionally, the similarity of many features in the spectra of
the macrocyclic NHC-based Cu(III) complexes and the
[Cu(CF3)4)]−species is consistent with a high valent Cu(III)
center in the latter. To corroborate this hypothesis, the Cu Kβ
Journal of the American Chemical Society
pubs.acs.org/JACS
Article
https://doi.org/10.1021/jacs.1c09505
J. Am. Chem. Soc. 2022, 144, 2520−2534
2524

\newpage

% Page 6

VtC spectra of these compounds were obtained and are
discussed below.
Cu Kβ XES Mainlines. As discussed in the introduction, X-
ray emission spectroscopies, particularly, Kβ mainlines, have
been shown to be diagnostic of the metal oxidation state and
spin state in other 3d transition metals. Variations in the Kβ
mainline features arise from metal 3p−3d exchange in the final
states.39 However, the utility of Kβ mainline emission
spectroscopy as a probe for the oxidation state of copper is
lacking. To initially explore the utility of Kβ mainlines for the
assignment of copper oxidation and spin states, the Kβ XES of
various copper halides were collected (Figure 3). The Kβ
mainlines of anhydrous CuCl, CuCl2, and CuF2 all exhibit very
similar Kβ1,3 energies (8904.6−8904.7 eV), with no significant
shifts between the various copper oxidation states. A slightly
more intense Kβ′ shoulder at 8894.7 eV is observed for the
CuCl2 sample compared to CuCl, as expected due to the
additional 3d hole for the CuII center. The energy splitting of
the Kβ1,3 and Kβ′ features (ΔE) is a marker of metal−ligand
covalency,35 as is also the intensity of the Kβ′ feature. For
CuF2, the energy of the Kβ′ feature increases while the
intensity decreases, yielding a diminished mainline splitting of
ΔE ≈8.2 eV, approximately 1.7 eV less than the splitting for
CuCl2 (Table S5). The variation of the splitting here
demonstrates that the covalency of the copper center has a
potentially significant influence on the appearance of the Kβ
mainline for samples of the same formal oxidation state.
However, as previously shown in Fe Kβ mainline studies, these
covalency interactions can also counteract the formal oxidation
and spin state effects, often making the interpretation of Kβ
mainlines difficult.35
The Kβ mainline XES spectra for [CuI/II/III(NHC2)]1/2/3+,
[CuIII(NHC4)]3+, and [Cu(CF3)4]−are also shown in Figure
3. Despite the clear variation in Cu oxidation state observed by
XAS, all [Cun+(NHC2)]n+ complexes exhibit almost super-
imposable Kβ mainlines and no shifts in the mainline Kβ1,3
peak’s energy (∼8904.9 eV), suggesting that no such analysis
can be conducted in this instance. As opposed to the maximum
splitting mainline differences observed for the copper halides,
the [Cun+(NHCx)]n+ series and [Cu(CF3)4]−complex all
exhibit similar ΔE values (Table S5) and a much smaller
splitting range (<0.6 eV). While it has often been shown that
Kβ mainlines are excellent reporters of electronic configuration
in ionic metal salts with lower dn electron counts, the lack of
differences for the reported Cu complexes is unsurprising given
their more complex electronic configuration and the generally
featureless Kβ mainline of more “ideal” copper halide salts.35
The remarkable similarities of these disparate formal Cu
oxidation states is reminiscent of the similar Kβ emission
spectra for covalent iron centers.41,55 This suggests that the
mainlines of these highly covalent molecular systems are poor
indicators of the Cu oxidation (and spin) state. In the next
section, the feasibility and sensitivity of Cu VtC XES as a probe
of copper oxidation state are investigated.
Cu Kβ VtC of Macrocyclic NHC Complexes. The VtC
spectrum of the Cu(I) species [CuI(NHC2)]+ is dominated by
a single, intense feature at 8975.9 eV, which increases in energy
to 8977.6 eV and finally to 8978.7 eV for [CuII(NHC2)]2+ and
[CuIII(NHC2)]3+, respectively (Figure 4). This suggests that
the transition energies in the VtC XES spectra are sensitive to
the oxidation state of the copper atom in this series where the
ligand environment remains consistent. The intensities of
features in the Kβ2,5 region of the emission spectrum have been
previously correlated to metal−ligand covalency.33,43,47 The
fitted VtC spectra show that the intensity of the Kβ2,5 region
predominantly arises from one high intensity emission band
(Figure S20); however, the intensity of this individual emission
band does not increase from Cu(I) to Cu(II) and then Cu(III)
to mirror the total emission intensity, which is likely due to
differences in molecular structure concomitant with oxidation
of the Cu center (vide infra).
In the Cu(II) and Cu(III) species, a shoulder grows into the
VtC spectrum on the low energy side of the main Kβ2,5 feature
(Figure S20) which is not observed in the Cu(I) speciesa
likely result of pyridine coordination and axial acetonitrile
ligands in the high valent species (vide infra).32 To higher
Figure 3. Cu Kβ XES mainlines. Kβ XES spectra have been
normalized by setting the integral of the mainline region to a value of
1.0 au.
Figure 4. Experimental (top) and calculated (bottom) VtC XES
spectra for the [Cun+(NHC2)]n+ series, [CuIII(NHC4)]3+, and
[Cu(CF3)4]−. Calculated spectra were shifted −5.1 eV and a 2.0 eV
(full-width half-maximum) Gaussian broadening was applied.
Journal of the American Chemical Society
pubs.acs.org/JACS
Article
https://doi.org/10.1021/jacs.1c09505
J. Am. Chem. Soc. 2022, 144, 2520−2534
2525

\newpage

% Page 7

energy of the main VtC features, between ∼8982−8983 eV,
there are emission features with noticeable intensity that are
not assigned as part of the Kβ2,5 feature. Overlays of the VtC
XES with the XAS of the [Cun+(NHC2)]n+ series (Figure S21)
show that these high energy VtC XES transitions overlap with
the weak pre-edge transitions of the XAS and are, thus, greater
than the Fermi energy. VtC XES and XAS in the single
electron transition picture probe the occupied and unoccupied
molecular orbitals, respectively. Hence in this simple picture,
no transition overlap is expected between the two spectros-
copies. However, features above the Fermi energy have also
been observed in Zn VtC XES and resonant VtC XES studies
of Cu(I) suggesting that these higher energy VtC XES peaks
likely arise from multielectron transitions.38,56
[CuIII(NHC4)]3+ exhibits an intense VtC peak at approx-
imately the same energy as the analogous [CuIII(NHC2)]3+
complex, albeit with much greater intensity (Figure 4). As in
the
[Cun+(NHC2)]n+
series,
the
spectrum
of
the
[CuIII(NHC4)]3+ species is dominated by a single intense
transition feature that is isoenergetic (8979 eV) with the
dominant feature of the [CuIII(NHC2)]3+ analogue (Figure
S20 and Table S6). The dramatic increase in intensity from
[CuIII(NHC2)]3+ to [CuIII(NHC4)]3+ is initially attributed to
the increased number of carbene donors present: two vs four.
NHC groups are significantly stronger sigma donors57−59 than
pyridyl nitrogens and thus result in more covalent metal−
ligand interactions and greater mixing of Cu np character into
the valence MOs.
In the [Cun+(NHC2)]n+ series, the Cu−C distances shorten
from Cu(I):1.94 Å to Cu(II):1.91 Å to Cu(III):1.88 Å, as
expected.32 In the Cu(I) species, the pyridines are essentially
noncoordinating, with long Cu−N distances of ∼2.6 Å;
however, in the Cu(II) and Cu(III) species the two pyridines
are coordinated, each with Cu−N distances of ∼2.16 and 1.97
Å, respectively (Figure S22). The coordination of the weaker
interacting pyridines is anticipated to bring about similarly
weak, although significant, contributions to the VtC emission
intensity in comparison to the carbene donors (vide infra).
Furthermore, inspection of the C−Cu−C angles in these
complexes shows that the angle in the Cu(I) species (173.9°)
and the Cu(III) species (172.7°) are closer to linear than that
in the Cu(II) species (169.0°). The more linear C−Cu−C σ
bonds in the Cu(I) and Cu(III) species will have better
overlap with the Cu npx orbitals, thus resulting in larger
intensity being imparted into the transition. This shows that
the decreasing Cu−C bond lengths in the Cu(II) species
relative to the Cu(I) species are counteracted by the geometry
changes resulting from the oxidation of the Cu center. Such an
angular dependence on the VtC spectrum of hydroxo and oxo-
bridged Fe dimers, as well as Ni-NOx and Cu-NOx (x = 1,2)
complexes, has been demonstrated previously.49,60 Hence, the
increase in the emission intensity observed in Cu(II) and
Cu(III) may be due to a combination of shorter Cu−
C(carbene) bond lengths, coordination of pyridine groups
resulting in significant overlap with the Cu npy orbitals, and
increased linearity in the CNHC−Cu−CNHC bond vector
facilitating further Cu npx mixing into the valence orbitals.60
It is clear from the experimental data of the [Cun+(NHC2)]n+
redox series that the VtC XES is sensitive to the oxidation state
of the Cu centers, manifesting as a stepwise increase in energy
of the dominant emission feature in the Kβ2,5 region.
To elucidate the origin of the changes in emission
intensities, DFT calculations were used to determine the
contributions of the ligand donor groups to the VtC spectra.
Ground-state DFT methods have been demonstrated to
accurately reproduce experimental XES spectra by calculating
the valence orbital →1s energy separations.33,39,45,47,51 Figure
4 shows the experimental and calculated VtC spectra for all five
complexes. For [CuIII(NHC4)]3+ and [Cu(CF3)4]−, the singlet
state was calculated to be significantly lower in energy than the
triplet state as expected for high valent copper, in agreement
with the experimental findings for [CuIII(NHC4)]3+ (vide
supra) and similar to the singlet spin state determined for
[CuIII(NHC2)]3+ (Table S7).14,17,32
The calculated VtC spectra for the [Cun+(NHC2)]n+ series
exhibit Kβ2,5 features with maxima at 8975.9, 8977.6, and
8978.7 eV for the Cu(I), Cu(II) and Cu(III) species,
respectively. The corresponding experimental maxima occur
at 8975.9 eV, 8977.2 and 8978.5 eV, showing good agreement
between experiment and theory (Figure S23). The DFT XES
calculations reproduce the increase in intensity of the Kβ2,5
feature as the oxidation state of the Cu center increasesin-
line with the experimental data shown in Figure S20.
Furthermore, DFT predicts the growth of a shoulder on the
low energy side of the main Kβ2,5 feature with increased
oxidation state, which is also observed experimentally as a
broadening to the lower energy side of the most intense VtC
feature (Figure S20). Inspection of the DFT-calculated Cu 1s
orbital energies shows that there is a stepwise contraction of
the core orbital energies with increased physical oxidation
state, decreasing by 2.7 eV from Cu(I) to Cu(II) and a further
2.1 eV from Cu(II) to Cu(III). This change in energy explains
the stepwise increase in the VtC Kβ2,5 feature, as the core
orbitals of the Cu center are more strongly influenced by the
increased Zeff of the high valent ion compared to the valence
orbitals. Hence, the energy gap between the MOs giving rise to
the dominant VtC transitions and the Cu 1s systematically
increases with physical oxidation state.
The single electron transition nature of the DFT VtC XES
calculations allows for individual transitions to be assigned to
the corresponding donor molecular orbitals.33,39 Figure 5
shows the calculated VtC spectra for the [Cun+(NHC2)]n+
series annotated with the donor MOs for the most intense
transitions of interest (additional MO assignments are shown
in Figure S24). The Cu(I) species is best described as a linear
two-coordinate molecule with local D2h symmetry, meaning
that the metal−ligand bonding interactions are comparatively
simplifiedthe intense Kβ2,5 VtC transitions arise from donor
MOs that are predominantly CNHC px with optimal overlap to
the acceptor Cu npx (Cu-L σ). The bonding MO of the
carbene presents predominant Cu-L σ overlap with the Cu dxz
and a minor π pz contribution (Figure S25). The large amount
of Cu npx orbital admixture yields the intense dipole allowed
feature, whereas the Cu-L π interaction does not present
favorable Cu pz overlap to produce appreciable dipole
intensity. The effect of this low-coordinate geometry for the
Cu(I) species on spectral intensity is also observed in the XAS,
exhibiting an intense pair of 1s →4p transitions at 8981.8 and
8984.8 eV in the pre-edge (Figure 2a)resulting from an
energetic splitting of the Cu px and py pairs as there is
effectively no ligation along the y-axis from the pyridyl
moieties. [CuII(NHC2)]2+ and [CuIII(NHC2)]3+ have coordi-
nating pyridine ligands and distort slightly from the higher
symmetry observed in the Cu(I) complex. Both the additional
ligands and geometrical distortion yield many more transitions,
which distributes the Cu np character across multiple final
Journal of the American Chemical Society
pubs.acs.org/JACS
Article
https://doi.org/10.1021/jacs.1c09505
J. Am. Chem. Soc. 2022, 144, 2520−2534
2526

\newpage

% Page 8

states;18 however, all dominant transitions exhibit Cu−CNHC σ-
character (Figure 5 and S10), as identified in [CuI(NHC2)]+.
This analysis is supported by the calculated polarized VtC
spectra, which show that for the [Cun+(NHC2)]n+ series,
almost all the intensity arises along the C−Cu−CNHC x vector,
with subtle contributions originating along the Npy−Cu−Npy y
vector for the Cu(II) and Cu(III) species, concomitant with
pyridine coordination (Figure S26), and in agreement with the
assertions made from the analysis of the experimental VtC fits
(Figure S20) and donor MOs calculated by DFT (Figure S27).
Conversely, [CuIII(NHC4)]3+ exhibits nearly equivalent VtC
intensity along both x and y CNHC−Cu−CNHC vectors as
expected.
Integration of the VtC spectra shows an ∼25\% increase in
intensity
upon
oxidation
of
[CuII(NHC2)]2+
to
[CuIII(NHC2)]3+, which is mostly due to the increase in
intensity of the main Kβ2,5 feature as a result of better overlap
between the CNHC 2px and Cu npx from a more linear CNHC−
Cu−CNHC bond vector and shorter Cu−CNHC distances, rather
than broadening or additional features, as shown by the change
in intensity of the dominant transition in the fits of the spectra
(Table S6 and Figure S20). Notably, there is a significant
increase in the amount of ligand np admixture going from the
Cu(II) (63.5\%) to the Cu(III) (70.0\%) species (Table S8).
The increase in the ligand np admixture in the valence MOs as
the Cu(II) species is oxidized is expected to lead to improved
ligand-mediated mixing of the Cu np character into the valence
MOs, imparting a larger amount of electric dipole allowed
character and, in turn, leading to increased emission intensity.
Furthermore, the geometric distortions caused by coordination
of the pyridine groups in the Cu(II) and Cu(III) species
clearly increases total metal−ligand interaction and enables
additional overlap of the ligand 2p and out-of-plane Cu npz
orbitals. Therefore, a mixture of shortening Cu−C bond
distances, additional Cu-pyridine interaction, and ligand-based
geometric distortions combine to explain the changes in
intensity observed between oxidation states for complexes in
the [Cun+(NHC2)]n+ series.
It is evident from the polarized VtC spectra that substitution
of the two pyridyl moieties in [CuIII(NHC2)]3+ with additional
NHCs in [CuIII(NHC4)]3+ results in a more intense VtC
spectrum. The tetra-NHC complex [CuIII(NHC4)]3+ exhibits
almost ideal local D4h symmetry in the CuC4 first coordination
sphere (Table S3), with the principal axis perpendicular to the
equatorial plane formed by the four carbene ligands (Scheme 2
and Figure S16). For this reason, almost all of the intensity in
the Kβ2,5 region stems from two degenerate donor MOs
formed by the C 2px and 2py like orbitals of adjacent NHC
moieties (Figure 5). Here, the mixing of Cu np character into
the valence MOs is achieved predominantly through CNHC−
CNHC σ bonding MOs which form an e doublet with very good
overlap with the Cu npx and npy orbitals (Figure 5).
The experimental VtC data shown here demonstrate that the
energy of the VtC emission is clearly sensitive to the oxidation
state at copper. Additionally, ground-state DFT methodologies
demonstrate the intensity of the emission is the product of
both the amount of metal−ligand covalency and the symmetry
of the complex ion. This series of complexes of the neutral
macrocyclic NHC-bearing ligands thus forms a convenient
calibration series for analysis of similarly covalent Cu species
such as [Cu(CF3)4]−.
Cu Kβ VtC of [Cu(CF3)4]−. The Cu Kβ VtC spectrum of
[Cu(CF3)4]−exhibits a predominant single feature to
significantly higher energy and intensity than both the
[CuI(NHC2)]+ and [CuII(NHC2)]2+ species, and ∼1.2 eV
higher in energy than both the Cu(III) NHC-based complexes
(Figure 4). The [Cu(CF3)4]−complex has an intensity that is
between the limits of [CuIII(NHC2)]3+ and [CuIII(NHC4)]3+
and also has a well resolved feature to lower energy at ∼8972
eV that is not observed in the Cu(III) NHCs. The fact that the
C−Cu−C angle (x-direction) in [Cu(CF3)4]−is more acute
(164.9°) than in [CuIII(NHC2)]3+ (172.7°), but exhibits more
intense VtC emission, suggests that the 4CF3
−donor set is
more covalent than the py2NHC2 set. However, it is clear that
per ligand, the CF3
−ligand produces less intensity than the
carbene ligand, after comparison to the VtC emission intensity
of the [CuIII(NHC4)]3+ complex. Typically, the amount of
metal np character that mixes into the valence MOs increases
with decreasing metal−ligand bond length. The average Cu−C
bond lengths in [Cu(CF3)4]−(1.96 Å) are significantly longer
than in [CuIII(NHC4)]3+ (1.88 Å), further indicating that the
CF3
−ligand is in fact a less covalent ligand and poorer σ-donor
than the NHC. The increase in energy of the Kβ2,5 feature in
[Cu(CF3)4]−relative to the NHC complexes is likely due to
the ionic nature of the CF3
−ligands, increasing the relative
energies of the donor MOs to the Cu 1s orbital as a result of
increased local negative charge. The differences between the
NHC and CF3
−ligand types is also expected to explain the
observation of a much more intense and well-resolved feature
to the low energy side of the dominant Kβ2,5 feature in
[Cu(CF3)4]−. Based on the experimental data, the VtC XES
supports the observations in the Cu K-edge XAS and supports
a high valent, highly covalent Cu(III) center in [Cu(CF3)4]−.
Figure 5. Calculated VtC spectra and selected transition donor MOs
for [CuI(NHC2)]+ (blue), [CuII(NHC2)]2+ (red), [CuIII(NHC2)]3+
(black), [CuIII(NHC4)]3+ (purple), and [Cu(CF3)4]−(light green).
MOs are plotted at an isosurface value of 0.05 au.
Journal of the American Chemical Society
pubs.acs.org/JACS
Article
https://doi.org/10.1021/jacs.1c09505
J. Am. Chem. Soc. 2022, 144, 2520−2534
2527

\newpage

% Page 9

DFT calculations were performed in order to elucidate the
nature of the VtC XES transitions and shed light on the origins
of the similarities and differences between the emission profiles
of the [Cu(CF3)4]−complex and those of the NHC-based
complexes. The higher symmetry and homoleptic C4 first
coordination sphere of the [CuIII(NHC4)]3+ (D4h) and
[Cu(CF3)4]−(D2d) complexes yields a narrower distribution
of final states for the donor MOs (Figure 5) than the
[CuIII(NHC2)]3+ complex. For [CuIII(NHC4)]3+ and [Cu-
(CF3)4]−, two effectively isoenergetic transitions from
degenerate orbitals of an e doublet (including both α and β
spin manifolds) are responsible for almost all the intensity of
the main Kβ2,5 feature (Figure 5). Despite the clear differences
of the neutral NHC and anionic CF3
−ligands, both complexes
exhibit remarkably similar valence MO structures. Analysis of
the atomic contributions to the donor MOs producing the
intense feature (Table S8) reveals that there is 9.2 and 10.1\%
Cu np admixture in [CuIII(NHC4)]3+ and [Cu(CF3)4]−,
respectively, further confirming the similarity in properties
between the two complexes. Additionally, this is significantly
larger than the ∼3.5−7.3\% calculated in the [Cun+(NHC2)]n+
series (Table S8).
The donor MOs in both the Cu(III) complexes with C4
ligand donor sets exhibit significant delocalization across the
first coordination sphere in the x- and y-directions, allowing for
a large degree of Cu npx,y mixing, as shown by the significant x-
and y-polarized intensity in Figure S26. As stated above, the
ionic nature of the CF3
−ligands increases the energies of the
donor MOs relative to the Cu 1s orbital: the dominant VtC
transitions in both the [CuIII(NHC4)]3+ and [Cu(CF3)4]−
complexes arise from HOMO−8/HOMO−9 in the former
and HOMO/HOMO−1 in the latter (Figure 5 and Table S4).
The fact that these highly similar donor MOs are frontier
orbitals in the [Cu(CF3)4]−complex, but slightly “buried” in
the [CuIII(NHC4)]3+ complex explains the higher energy Kβ2,5
feature in the [Cu(CF3)4]−complex relative to the NHC-
bearing complexes.
In the experimental spectrum of [Cu(CF3)4]−, a unique and
sharp lower energy feature at ∼8972 eV is well resolved
(Figure 4). DFT calculations show that this feature arises
almost exclusively from π-interactions between fluorine 2px,y
and carbon 2px,y of the CF3
−ligands, which are ideally
positioned to form σ-type interactions with the Cu atom along
the x- and y-directions (Figure 5). These MOs exhibit
significant overlap with the Cu npx,y orbitals, resulting in
∼3.1\% Cu p character and ∼88.8\% total ligand p character,
which ultimately explains the increased intensity in this region
compared to the NHC-based complexes, which exhibit closer
to ∼1.0\% Cu np admixture into the donor MOs in this region.
The great similarities between the both the XAS and VtC
XES of the [CuIII(NHC4)]3+ and [Cu(CF3)4]−complexes,
coupled with theory, suggest a similar electronic structure
across both complexes. Furthermore, the [Cun+(NHC2)]n+
series and [CuIII(NHC4)]3+ complexes exhibit clear sensitivity
to the copper’s physical oxidation state, where the observed
shifts in energy of the main Kβ2,5 features are a function of the
deepening of the Cu 1s orbital through increased Zeff.
■DISCUSSION
Relationship between Cu 1s Ionization Energy and
Oxidation State. For the [Cun+(NHC2)]n+ series, the VtC
data presented in the current study are in-line with the
previously reported Cu K-edge XAS data, exhibiting clear
spectroscopic energy shifts correlated to the formally assigned
oxidation state.32 These trends are mirrored by the Cu VtC
XES spectra of the [Cun+(NHC2)]n+ series, where the
dominant Kβ2,5 feature increases incrementally in energy
with increasing physical oxidation state (Figure 4). Having
demonstrated that the XAS pre/rising edge features and VtC
XES data are accurately reproduced by (TD)DFT calculations,
we can use the calculated electronic structures to better
Figure 6. Qualitative representation of the electronic structures of the complexes discussed in the present analysis. Dashed orange lines correspond
to the IWAE of the quadrupole-only VtC spectrum and represent the approximate energetic center of the Cu 3d manifold. Solid pink lines
represent the HOMO energy determined from the DFT-calculated Cu Kβ VtC spectra. Solid blue lines represent the LUMO energy as determined
experimentally from the Cu K-edge XAS. The relative energies of the Cu 1s and HOMO for each complex are shown in black and were determined
from DFT calculations (B3LYP/ZORA-def2-TZVP/ZORA), whereas the relative energy of the LUMO shown in green was determined from
TDDFT calculations (B3LYP/ZORA-def2-TZVP/ZORA).
Journal of the American Chemical Society
pubs.acs.org/JACS
Article
https://doi.org/10.1021/jacs.1c09505
J. Am. Chem. Soc. 2022, 144, 2520−2534
2528

\newpage

% Page 10

understand the origins of the energy shifts in the XAS and VtC
spectra. This combined XAS/XES approach enables the
formation of a more complete electronic structure diagram
(Figure 6) and correlation of the spectroscopic trends to the
relative energies of the core Cu 1s orbital and the valence
MOs.
While [CuIII(NHC2)]3+ and [CuIII(NHC4)]3+ are both
assigned as Cu(III) centers, their pre-edge energies differ by
approximately 1 eV. Inspection of the relative orbital energies
shows that the Cu 1s orbital of the [CuIII(NHC4)]3+ species is
∼0.5 eV higher in energy than in the [CuIII(NHC2)]3+
analogue, however, it is still nearly 2 eV lower than the Cu(II)
1s orbital energy of [CuII(NHC2)]2+. The LUMO of
[CuIII(NHC4)]3+ is shifted ∼1.5 eV higher in energy than
the [CuIII(NHC2)]3+ species; the combined 1s and LUMO
energy difference yield the ∼1 eV higher energy shift of the
[CuIII(NHC4)]3+
pre-edge
transition
compared
to
[CuIII(NHC2)]3+. This demonstrates how the more covalent
interactions of the NHC4 vs py2NHC2 first coordination
spheres with the Cu center modulates the pre-edge transition
energy as a function of the increased relative energy of the
3dx
2−y
2−2px,y σ* antibonding orbital (Figure S18), despite
minimally raising the energy of the Cu 1s orbital.
The [Cu(CF3)4]−complex exhibits a similar XAS spectrum
to that of both [CuIII(NHC2)]3+ and [CuIII(NHC4)]3+ (Figure
2), with a Cu 1s orbital energy calculated to be significantly
more negative than that of the established Cu(II) center in
[CuII(NHC2)]2+. The observed energy of the pre-edge feature
in
the
XAS
of
[Cu(CF3)4]−
is
intermediate
to
[CuIII(NHC2)]3+ and [CuIII(NHC4)]3+, which may be
interpreted as increased metal−ligand covalency across the
ligand series such that py2NHC2 < 4CF3
−< NHC4. The
modulation of the observed transition energies is in-line with
the anticipated increase in ligand field resulting from these
strong σ donating ligands. This trend is also observed in the
VtC emission intensities, which appear to increase with
increased metal−ligand covalency (Figure 5).
As mentioned above, (TD)DFT methods show that the XAS
and VtC XES spectroscopic trends in the [Cun+(NHC2)]n+
series are due to the contraction of the Cu 1s orbital to deeper
binding energy upon oxidation of the central Cu ion: the Cu 1s
orbital energy decreases by 2.8 eV between the Cu(I) and
Cu(II) species, and by a further 2.5 eV between the Cu(II) and
Cu(III) species (Figure 6). The significant contraction of the
core Cu 1s orbitals compared to the more subtle contraction of
the frontier/valence MOs explains the increasing energy of the
Cu K-edge pre-edges and VtC emission energies of the Kβ2,5
features for the [Cun+(NHC2)]n+ series. Similarly, the Cu 1s
binding energy in both [CuIII(NHC4)]3+ and [Cu(CF3)4]−is
calculated to be to 4.9 and 3.7 eV lower in energy than in
[CuI(NHC2)]+, respectively, supporting the notion that the
spectroscopic features and their relative energies are due to the
increase in Cu Zeff and its powerful capacity to reflect the
copper ion’s formal oxidation state. These results strongly
support a physical Cu(III) oxidation state and low-spin 3d8
electronic configuration for the [Cu(CF3)4]−complex, even
with the strong charge donation from the anionic CF3
−ligands.
The [CuIII(NHC2)]3+ and [CuIII(NHC4)]3+ complexes have
Cu 1s orbitals that are 1.7 and 1.2 eV lower in energy than that
of the [Cu(CF3)4]−complex, respectively, which is likely due
to the fact that the neutral macrocyclic bis-NHC and tetra-
NHC ligands provide less charge donation to the high valent
Cu centers than the CF3
−ligands.
This combined XAS/VtC XES approach allows for
evaluation of the relative energies of the LUMO and valence
MOs as a function of the Zeff of the Cu, as demonstrated in
Figure 6, which emphasizes the impact of the physical
oxidation state on the core X-ray spectroscopies employed in
this study. While the 3d orbitals do not contribute significant
intensity to the VtC XES, the average energy of this manifold
may be extracted from the quadrupole (3d →1s) contribution
of the calculated VtC XES spectra (Figure S28). For the Cu(I)
species [CuI(NHC2)]+, the 3d manifold is represented by the
five highest occupied MOs (Figure 6 and Figure S24), which
are easily identified. However, in the Cu(II) and Cu(III)
species the increased metal−ligand covalency and coordination
number causes the highest valence MOs to have significantly
more ligand character, resulting in MOs that do not reflect the
localized geometries of the spherical harmonics that are
typically associated with “pure” 3d orbitals. However, analysis
of the quadrupole contribution to the VtC XES effectively
reports the distribution of the 3d orbital character without the
need for MO analysis. The effective VtC transition energy of
the 3d manifold is plotted in Figure 6, and when viewed on the
absolute energy scale, it is clear that as oxidation state increases
across the [Cun+(NHC2)]n+ series, the 3d manifold and the 1s
core orbital are found to exhibit deeper binding energy.
Accompanying the decreasing 3d manifold energy with Cu 1s
orbital contraction is a relative expansion of the 3d manifold,
where the energy range of the observed quadrupole transitions
increases, and is most obvious to higher energies (Figure S28).
The dominant VtC XES transition observed for these
organocopper complexes arises from MOs to the higher
energy edge of the d-manifold. A consistent picture develops
from the [Cun+(NHC2)]n+ series, where the increased
dispersion of d-orbital character across the valence MOs is
brought about by an increased Cu physical oxidation state and
concomitant changes in metal−ligand covalency. Considering
the [CuIII(NHC4)]3+ complex, the quadrupole-only VtC
spectrum is highly similar to the [CuIII(NHC2)]3+ analogue,
and noticeably different from the Cu(I) and Cu(II) species.
The [CuIII(NHC2/4)]3+ complexes exhibit quadrupolar inten-
sity at higher energies than in the Cu(I) or Cu(II) species, and
the dominant feature in the spectrum is less well-resolvedin-
line with a distribution of the d-orbital character across a larger
number of MOs. These observations support the trends in the
experimental VtC XES and XAS, and by using this combined
experimental approach, these complexes form a useful
calibration series for the determination of both physical
oxidation state and metal−ligand covalency in the
[CuIII(NHC4)]3+ and [Cu(CF3)4]−complexes. Given the
aforementioned calibrations, we can now rationalize both the
Cu K-edge XAS and VtC XES of the [Cu(CF3)4]−complex in
terms of its physical oxidation state and proposed electronic
configurations.
Physical Oxidation State of [Cu(CF3)4]−. The Cu K-edge
XAS, VtC XES, and the accompanying DFT calculations
provide significantly strong evidence for a low-spin d8 Cu(III)
configuration on the Cu center of the [Cu(CF3)4]−complex.
This conclusion is in contradiction to recent studies of this
complex, where it has been suggested that an “inverted ligand
field” dictates that the MOs with majority Cu 3d character are
fully populated, and a d10 electronic configuration is the most
accurate description of the electronic structure (vide infra).10
The “inverted ligand field” phenomenon in Cu complexes
was discussed in much further detail in a 2016 review by
Journal of the American Chemical Society
pubs.acs.org/JACS
Article
https://doi.org/10.1021/jacs.1c09505
J. Am. Chem. Soc. 2022, 144, 2520−2534
2529

\newpage

% Page 11

Hoffmann et al.,17 and in the pursuit of transparency, the
authors stated that there remained no consensus as to what
inverted ligand fields imply for the electronic structure in Cu
complexes. We are of the opinion that increased covalency of
the ligand and concomitant increase of ligand character in the
valence MOs does not necessarily dictate the presence of an
inverted ligand field. In classic molecular orbital theory, the
energy difference between the d-orbitals and the ligand is used
to describe the covalent character, where large separations are
termed ionic and those approaching equal energy are highly
covalent. However, the suggestion that the lowering of the d-
orbital character in the LUMO below 50\% yields an inverted
ligand field is not a unique suggestion, only a renaming, and
the covalency of the system would remain equal for d orbital/
ligand orbital displacements equidistant from the 50\% mark.
For example, this exact case has been more simply described as
stabilization (or destabilization) of the metal center through
“inverted bonding”,26 as the ligand field of the metal center
itself is not inverted.1,61 In metal complexes with such large
amounts of metal−ligand covalency, referring to select few
MOs with slightly larger Cu 3d admixtures as the Cu d-
manifold is problematic, as the distribution of Cu character
among the valence orbitals is vast and significantly lower than
that which is seen in Cu(I) systems.
In 2016, Lancaster and co-workers employed a combination
of Cu K-edge XAS and Cu 1s−2p resonant inelastic X-ray
scattering (RIXS) in order to elucidate the electronic structure
of the [Cu(CF3)4]−complex.14 Significant attention was
focused toward the 32\% Cu 3d contribution to the LUMO
(Cu 3dx
2−y
2−L σ*, b2 in D2d symmetry), leading to the
suggestion that the pre-edge transitions in the Cu K-edge XAS
spectrum should be described as excitations of the Cu 1s core
electrons to predominantly ligand MOs rather than to Cu 3d
orbitals, implying a 3d10 electronic configuration and
“inversion” of the ligand field on Cu.14 It would later be
suggested that all Cu(III) systems exhibit inverted ligand fields
and should be described physically as d10 systems, with the
exception of the [CuF6]3−species.10
In 2015, Tomson et al. investigated the Cu K-edges of a
range of covalent β-diketiminate Cu complexes with co-ligands
exhibiting various degrees of π acidity.18 In this work, pre-edge
features with energies typical of Cu(III) species were found to
be accompanied by uncharacteristically low energy rising
edges, specifically Cu 1s →4p excitations, something which
was also subsequently observed in bipyridine coordinated
Cu(I/II) complexes.62 Here the authors stress that an analysis
based solely on the pre-edge energies of these complexes
would have led to a Cu(III) assignment on many complexes
that were shown via quantum chemical calculations to have
ground states dominated by Cu(I)/Cu(II) configurations.
Hence, this study emphasizes the importance of examining
both the pre-edge and rising edge energy trends in the
evaluation of the electronic structure of copper complexes.
Unlike the low valent copper complexes studied by Tomson et
al.,18 which exhibited higher energy pre-edges arising from Cu
→π* transitions and low energy rising edge energies, the
formally Cu(III) complexes investigated in the present study
exhibit both higher energy pre-edge features (relative to the
Cu(I) and Cu(II) species) and also higher energy rising edge
features. Finally, this trend is extended to the VtC XES, where
the Cu(III) systems exhibit higher energy Kβ2,5 features
relative to the Cu(I)/(II) systems which supports a larger Zeff
on the three formally Cu(III) complexes. This reinforces the
need to employ a more holistic approach to the investigation
and evaluation of the electronic structures of highly covalent
Cu complexes. It is here where the combination of the Cu K-
edge XAS and Cu Kβ VtC XES can be combined to overcome
limitations in assignment of oxidation state from XAS alone.
The stepwise increase in VtC emission energy with increased
formal oxidation state in the [Cun+(NHC2)]n+ series has
displayed a clear sensitivity and utility of the VtC XES
technique to the physical oxidation state of Cu. In combination
with DFT calculations, the parent valence transitions of this
series reflect the energy of the core 1s orbital and the effective
nuclear charge (i.e., oxidation state). DFT calculations also
identified the remarkably similar electronic structures of the
[Cu(CF3)4]−and [CuIII(NHC4)]3+ complexes: both com-
plexes exhibit Cu 1s orbital ionization energies entirely
consistent with that observed for [CuIII(NHC2)]3+, while
also displaying strikingly similar valence MOs to one another,
leading to their intense and high energy Kβ2,5 emission
features. By simple analogy, one is led to conclude that
[Cu(CF3)4]−is in fact best described as a d8 Cu(III) center
and that the d10 Cu(I) oxidation is not supported.
Importance of Observables in Physical Oxidation
State Assignment. The [CuCl4]2−ion is often used as a
reference point for ligand field inversion, due to the extensive
amount of experimental data available. The spin population on
Cu in the D4h [CuCl4]2−ion has been experimentally
determined as 56−67\% via a range of methods.63 Thus, it is
accepted that the resulting ligand field is conventional, with the
relative energies of the Cu frontier orbitals greater than those
of the ligand(s). In order to computationally reproduce a
>50\% Cu d-orbital admixture into the HOMO of the
[CuCl4]2−complex via DFT methods, the amount of
Hartree−Fock (HF) exchange that is employed by the selected
functional must be calibrated.63,64 Analysis of the LUMOs (i.e.,
the α-spin lowest unoccupied MO in a spin-unrestricted
description) of the Cu(III) systems investigated in the present
study shows that there is diminished Cu nd admixture (Table
S10), dropping below the 50\% threshold frequently employed
as the cutofflimit for this select definition of an “inverted
ligand field”. However, this is also the case for the HOMO
(i.e., the α-spin highest occupied MO in a spin-unrestricted
description) of the Cu(II) center in [CuCl4]2−analyzed by
both Löwdin and Mulliken population analysis (B3LYP/def2-
TZVP). In the [Cun+(NHC2)]n+ redox series, the Cu(I)
species has a HOMO with 58.3\% Cu nd admixture (Löwdin),
decreasing to 29.5\% in the HOMO of the Cu(II) species, and
then increasing again to 42.0\% in the LUMO of the Cu(III)
species. These MOs are effectively the analogous Cu 3dx
2−y
2−
2px,y σ* antibonding orbitals across the redox series.
Despite the <50\% Cu 3d admixture into the HOMO in
[CuII(NHC2)]2+, the Löwdin spin population on the Cu atom
is ∼61.5\%. Similar observations are made when observing the
Löwdin population of the HOMO of D4h [CuCl4]2−, which is
calculated to have only 20.2\% Cu nd admixture using the same
B3LYP/def2-TZVP method. Comparatively, the calculated
Löwdin population of the unrestricted natural orbitals (UNOs)
gave a Cu nd admixture that was much closer to the calculated
spin population (Table S10). This is also true for the Cu(II)
complex [CuII(NHC2)]2+, which is calculated to have a Cu
Löwdin spin population of 61.5\% and 61.0\% Cu nd admixture
in the HOMO from Löwdin population analysis of the UNOs.
This is in stark contrast to the 29.5\% Cu nd admixture in the
HOMO according to the Löwdin population analysis of the
Journal of the American Chemical Society
pubs.acs.org/JACS
Article
https://doi.org/10.1021/jacs.1c09505
J. Am. Chem. Soc. 2022, 144, 2520−2534
2530

\newpage

% Page 12

canonical orbitals, which would qualify as an “inverted ligand
field”, despite experimental EPR evidence for a significantly
large Cu hyperfine and thus Cu admixture into the ground
state wave function.32 In the case of the closed-shell Cu(I) and
Cu(III) complexes, Cu nd admixtures in the HOMO and
LUMO, respectively, according to the population analyses on
the UNOs, are consistently much less than 10\%, which is
unrealistic in the case of the Cu(I) complex [CuI(NHC2)]+
considering the HOMO has dominant Cu 3dx
2−y
2 character
(Figure S24). Inspection of the quasi-restricted orbitals
(QROs) shows that both in the Löwdin and Mulliken
population analyses, the Cu nd admixture into the HOMO/
LUMO for the closed shell Cu(I)/Cu(III) complexes closely
match that of the canonical orbitals, whereas in the open shell
Cu(II) complexes, the QROs are a much closer match to the
UNOs than they are to the canonical orbitals. The wide
variation in the degree of Cu admixtures suggests that
population analyses of the MOs should be cautiously
implemented when defining the electronic structures of Cu
complexes, especially as there may be large disparities between
the spectroscopically observable trends and the mathematically
constructed MO populations.
The discrepancies between the values produced by various
population analyses clearly highlight the crucial role of
assessing the spectroscopic signatures from core X-ray
spectroscopies when assigning oxidation states in high valent
organocopper species. The Cu(III) assignment for [Cu-
(CF3)4]−is entirely consistent among all of the spectroscopic
observables for this complex to date. Here, it is demonstrated
how the observed transitions for the Cu K-edge XAS and VtC
XES of [Cu(CF3)4]−are due to a large Zeff, consistent with a
Cu(III) oxidation state.
The low intensity of the Cu L-edges in [Cu(CF3)4]−,
derived from 1s−2p RIXS experiments, have been interpreted
as a diminished d-hole character, leading to invocation of an
inverted ligand field and a Cu(I) d10 configuration. However,
the energies of the observed L-edge transitions in [Cu(CF3)4]−
are in agreement with numerous other formally assigned
Cu(III) complexes.9,10,28 While the L-edges of [Cu(CF3)4]−
and other Cu(III) centers do indeed exhibit decreased
intensities per d-hole, this is clearly correlated to the increased
ligand covalency, as probed by ligand K-edge XAS.28 This
interpretation of the low Cu d-character measured by the L-
edge XAS is also consistent with S = 1/2 Cu(II) thiolate
centers such as plastocyanin, which exhibits typical blue-copper
EPR signatures and ∼41\% Cu 3d admixture into the ground
state wave function,31 being described as highly covalent
copper-based paramagnetic centers with large amounts of
ligand character contributing to the SOMO, as opposed to
Cu(I) systems with ligand-based radical(s), as would be the
required description by employing this interpretation of an
inverted ligand field.
Furthermore, the lack of classic d−d transitions in the UV−
vis spectrum of [Cu(CF3)4]−is expected for the large ligand
field expansion in the presence of the high valent copper ion,
and its increased Zeff (Figure 6). From the combination of VtC
XES and XAS studies as well as computational investigations, it
is demonstrated that the MOs responsible for the dominant
VtC feature and pre-edge transition in [Cu(CF3)4]−are the
ligand-based HOMO/HOMO−1 and the Cu 3dx
2−y
2 localized
LUMO, respectively. Therefore, the combination of VtC XES
and Cu K-edge XAS is a direct measure of the HOMO−
LUMO gap for [Cu(CF3)4]−. Electronic excitations between
the HOMO and LUMO are calculated by TDDFT (Figure
S29) to have negligible oscillator strength and are thus not
observed in the reported UV−vis spectrum.14 Due to the large
ligand field, the lowest energy d−d transition with non-
negligible oscillator strength is calculated at ∼36,600 cm−1 and
is one with contributions from two valence MOs with
significant Cu 3dz
2 character, forming a dz
2 →dx
2−y
2
(LUMO) transition (Figure S29). A pair of much more
intense transitions occur within the same absorption band,
stemming from a combination of excitations from different
donor orbitals into the LUMO: the ligand based e doublet
(C−C σ bonding, also responsible for the VtC XES intensity)
and the 3dxz/yz e doublet. To even higher energy is a yet more
intense absorption band involving identical orbital sets. The
involvement of the C−C σ-bonded ligand MOs in these
electronic excitations provides significant charge transfer
character to the transitions; thus, the intensity is much larger
than other d−d transitions mentioned (Figure S29). A
remarkably similar pair of intense absorption bands is observed
in the [CuIII(NHC4)]3+ complex (Figure S30), as expected
from the similarities in electronic structures. The highest
energy absorption band shown in Figure S30 is isoenergetic
between both the [Cu(CF3)4]−and [CuIII(NHC4)]3+ com-
plexes, and stems from highly similar donor MO sets into the
largely 3dx
2−y
2 LUMO. This demonstrates how the high valent
and highly covalent Cu(III) complexes with C4 donor sets
exhibit atypical high energy and high intensity pseudo d−d
transitions due to the increased charge transfer character that is
imparted by the involvement of σ-type MOs formed by
adjacent carbon atoms in the first coordination sphere.
■CONCLUSIONS
A combination of Cu K-edge XAS and Cu Kβ VtC XES has
been used to spectroscopically probe the oxidation state in a
suite of organometallic Cu complexes and has provided strong
evidence for the presence of a low-spin d8 Cu(III) center in
three separate ligand systems. The pre-edge, edge, and white
lines of the complexes studied exhibit stepwise increases in
energy as the physical oxidation state of the Cu center is
increased, which is supported by TDDFT calculations.
Similarly, both the energies and intensities of the dominant
Kβ2,5 features in the VtC XES increase in a predictable fashion
upon Cu oxidation and are well reproduced by DFT
calculations. The most intense VtC features occur in the
Cu(III) complexes, with the more symmetric pseudo-D4h
[CuIII(NHC4)]3+ and D2d [Cu(CF3)4]−species exhibiting the
largest intensitythe result of increased Cu np character being
imparted into the valence MOs in these highly covalent C4
coordination environments. The Kβ2,5 intensity therefore
clearly scales with the σ-donor strength of the ligands;
however, there is also an angular dependence on this intensity
in the macrocyclic NHC-based complexes, where the intensity
increases as the trans C−Cu−C angles approach linearity.
Analysis of the valence MOs shows that the Cu
contributions to the antibonding 3dx
2−y
2-L orbital in both
Cu(II) (HOMO) and Cu(III) species (LUMO) is significantly
less than the 50\% threshold that is often used to describe the
limit between regular and inverted ligand fields. Analysis of the
HOMO in the ubiquitous [CuCl4]2−standard shows that
despite having ∼60\% spin population on Cu atom, the
calculated Cu 3d admixture in the highest occupied α spin MO
is lower than all of the complexes investigated in the present
study. Although MO population analysis is demonstrated to be
Journal of the American Chemical Society
pubs.acs.org/JACS
Article
https://doi.org/10.1021/jacs.1c09505
J. Am. Chem. Soc. 2022, 144, 2520−2534
2531

\newpage

% Page 13

able to explain the observable trends in the experimental XAS
and VtC XES data, the values extracted from this mathematical
procedure can vary dramatically based on the type of orbital
localization method and should therefore be used qualitatively
rather than quantitatively.65,66 The results of the XAS and VtC
XES studies, in combination with (TD)DFT calculations, are
in strong support of assignment of the formally Cu(III)
systems as low-spin d8 ions, characterized by highly covalent
bonding networks and high energy pre-edge and rising edge
Cu K-edge XAS and VtC features.
In summary, we have shown that by combining Cu K-edge
XAS, VtC XES and computational methods, a more complete
description of the electronic structure of high-valent copper
complexes is achieved. Herein, the controversial [Cu(CF3)4]−
ion is shown to be best described as a low-spin d8 system with
a dx
2−y
2 LUMO that is engaged in highly covalent bonding
interactions with the homoleptic (CF3
−)4 ligand donor set.
Finally, we emphasize the need to take a comprehensive
approach when defining the electronic structure of highly
covalent copper complexes. Careful evaluation of the
spectroscopic transition energies and intensities, preferably
utilizing more than one experimental method, is needed to
arrive at a more comprehensive understanding
■ASSOCIATED CONTENT
*
sı Supporting Information
The Supporting Information is available free of charge at
https://pubs.acs.org/doi/10.1021/jacs.1c09505.
Materials and Methods; analytical characterization of
[CuIII(NHC4)](PF6)3, additional XAS and XES, DFT
computations, xyz coordinates (PDF)
Accession Codes
CCDC 2104793 contains the supplementary crystallographic
data for this paper. These data can be obtained free of charge
via www.ccdc.cam.ac.uk/data\_request/cif, or by emailing
data\_request@ccdc.cam.ac.uk, or by contacting The Cam-
bridge Crystallographic Data Centre, 12 Union Road,
Cambridge CB2 1EZ, UK; fax: +44 1223 336033.
■AUTHOR INFORMATION
Corresponding Author
George E. Cutsail III −Max Planck Institute for Chemical
Energy Conversion, 45470 Mülheim an der Ruhr, Germany;
Institute of Inorganic Chemistry, University of Duisburg-
Essen, 45117 Essen, Germany;
orcid.org/0000-0002-
7378-9474; Email: george.cutsail@cec.mpg.de
Authors
Blaise L. Geoghegan −Max Planck Institute for Chemical
Energy Conversion, 45470 Mülheim an der Ruhr, Germany;
Institute of Inorganic Chemistry, University of Duisburg-
Essen, 45117 Essen, Germany;
orcid.org/0000-0001-
9543-102X
Yang Liu −Institute of Inorganic Chemistry, University of
Göttingen, 37077 Göttingen, Germany;
orcid.org/0000-
0002-2486-8520
Sergey Peredkov −Max Planck Institute for Chemical Energy
Conversion, 45470 Mülheim an der Ruhr, Germany
Sebastian Dechert −Institute of Inorganic Chemistry,
University of Göttingen, 37077 Göttingen, Germany;
orcid.org/0000-0002-2864-2449
Franc Meyer −Institute of Inorganic Chemistry, University of
Göttingen, 37077 Göttingen, Germany;
orcid.org/0000-
0002-8613-7862
Serena DeBeer −Max Planck Institute for Chemical Energy
Conversion, 45470 Mülheim an der Ruhr, Germany;
orcid.org/0000-0002-5196-3400
Complete contact information is available at:
https://pubs.acs.org/10.1021/jacs.1c09505
Funding
APC Funding Statement: Open access funded by Max Planck
Society.
Notes
The authors declare no competing financial interest.
■ACKNOWLEDGMENTS
We thank Dr. Sandeep Gupta (University of Göttingen) for
collecting and analyzing the magnetic susceptibility data of
[CuIII(NHC4)](PF6)3. Y.L. is grateful to the Alexander von
Humboldt Foundation for a postdoctoral fellowship. G.C.,
B.L., S.P., and S. DeBeer. acknowledge financial support from
the Max Planck Society. S. DeBeer. also acknowledges funding
from the European Research Council under the European
Union’s Horizon 2020 research and innovation program
(Grant Agreement 856446). The PINK Beamline is supported
through the Max Planck Society, Helmholtz-Zentrum Berlin,
and the Energy Materials In-Situ Laboratory Berlin (EMIL).
SSRL is acknowledged for allocation of beamtime (BL6-2),
and Thomas Kroll is thanked for beamline assistance. SSRL is
supported by the U.S. Department of Energy, Office of
Science, Office of Basic Energy Sciences under Contract No.
DE-AC02-76SF00515. We acknowledge SOLEIL for provision
of synchrotron radiation facilities, and we would like to thank
Gautier Landrot for assistance in using the SAMBA beamline.
This project has been partly funded by the Deutsche
Forschungsgemeinschaft (DFG) Project Number 423268549
(INST 186/1327-1 FUGG) and the Nds. Ministerium für
Wissenschaft und Kultur (MWK).
■REFERENCES
(1) Solomon, E. I.; Heppner, D. E.; Johnston, E. M.; Ginsbach, J. W.;
Cirera, J.; Qayyum, M.; Kieber-Emmons, M. T.; Kjaergaard, C. H.;
Hadt, R. G.; Tian, L. Copper Active Sites in Biology. Chem. Rev. 2014,
114 (7), 3659−3853.
(2) Phipps, R. J.; Gaunt, M. J. A Meta-Selective Copper-Catalyzed
C-H Bond Arylation. Science 2009, 323 (5921), 1593−1597.
(3) Chen, B.; Hou, X. L.; Li, Y. X.; Wu, Y. D. Mechanistic
Understanding of the Unexpected Meta Selectivity in Copper-
Catalyzed Anilide C-H Bond Arylation. J. Am. Chem. Soc. 2011, 133
(20), 7668−7671.
(4) Le, C.; Chen, T. Q.; Liang, T.; Zhang, P.; Macmillan, D. W. C. A
Radical Approach to the Copper Oxidative Addition Problem:
Trifluoromethylation of Bromoarenes. Science 2018, 360 (6392),
1010−1014.
(5) Hossain, A.; Bhattacharyya, A.; Reiser, O. Copper’s Rapid Ascent
in Visible-Light Photoredox Catalysis. Science 2019, 364 (6439).
DOI: 10.1126/science.aav9713.
(6) Kainz, Q. M.; Matier, C. D.; Bartoszewiez, A.; Zultanski, S. L.;
Peters, J. C.; Fu, G. C. Asymmetric Copper-Catalyzed C-N Cross-
Couplings Induced by Visible Light. Science 2016, 351 (6274), 681−
685.
(7) Liu, H.; Shen, Q. Well-Defined Organometallic Copper(III)
Complexes: Preparation, Characterization and Reactivity. Coord.
Chem. Rev. 2021, 442, 213923.
Journal of the American Chemical Society
pubs.acs.org/JACS
Article
https://doi.org/10.1021/jacs.1c09505
J. Am. Chem. Soc. 2022, 144, 2520−2534
2532

\newpage






--- Page 34 ---
 
34 
XYZ Coordinates for Orca Calculations 
[CuI(NHC2)]+ 
Charge = 1+, Spin Multiplicity = 1 
  Cu  -0.00740413299260     -0.04079336237917      0.13207803040467 
  N   -0.02306510512784     -2.62451096464465      0.39810987921352 
  N   0.00932859614757      2.56090157725132      0.28689515106662 
  C   1.94728414746307     -0.07464419591113      0.22304221810076 
  C   -1.96070678328296      0.00365151709803      0.24885425026110 
  C   1.08157935414333     -2.80336943362877      1.12738717334633 
  C   -1.21411673613655     -2.79297614879768      0.97584022771510 
  C   -1.08332126019588      2.76462195233140      1.02745388054978 
  C   1.20918221142673      2.74988236796822      0.83932751055978 
  N   2.77476254700957     -1.13712084620435      0.39209544958973 
  N   2.79304105107554      0.98787826466697      0.18784269586343 
  N   -2.78672815754387      1.07280132298858      0.37910296228816 
  N   -2.80684976197257     -1.05915199921355      0.26438925445634 
  C   2.37810980747925     -2.55535313027989      0.40701562104720 
  C   1.03649430977973     -3.16881219046406      2.46758777732075 
  C   -1.35064324389426     -3.17142403630490      2.30676528152688 
  C   -2.40085380900937     -2.46631470888475      0.11221353571653 
  C   -2.39162657503318      2.49119399335755      0.33812444381695 
  C   -1.01655894519240      3.17624761781704      2.35313538783983 
  C   1.36735837672303      3.17618543289052      2.15318018745124 
  C   2.38373841133867      2.38632893324904     -0.02548123351970 
  C   4.10057715692895     -0.74861859514705      0.44610284290877 
  C   4.11180015935928      0.59674217931508      0.31921745065780 
  C   -4.11228945727196      0.68712224807323      0.45745623321823 
  C   -4.12467488342361     -0.66230320389551      0.38682817435480 
  C   -0.20365214973805     -3.35912182510847      3.06403240449292 
  C   0.23298074223064      3.38932364419960      2.92192097796087 
  H   3.18870313031190     -3.10199987409391      0.88468436480272 
  H   2.28322449815311     -2.90045117229094     -0.62234805289014 
  H   1.94963117983410     -3.31241259409035      3.02998583768353 
  H   -2.33170808456379     -3.31416484129875      2.74082452252246 
  H   -2.14975535790919     -2.65555525230473     -0.93040014891398 
  H   -3.26187618230452     -3.07406076961440      0.38354869009515 
  H   -3.19455076072294      3.05305206474972      0.81098475814832 
  H   -2.31558271709328      2.80123440237509     -0.70399280345024 
  H   -1.92058719760087      3.33742714306229      2.92528875481498 
  H   2.35541366793588      3.33417698636880      2.56562933115066 
  H   2.11753966469871      2.52565130381315     -1.07220997868685 
  H   3.24760556566286      3.00677083226666      0.20491660408108 


--- Page 35 ---
 
35 
  H   4.90492510117534     -1.45075645416561      0.57301593914522 
  H   4.92890452820429      1.29540084567437      0.31498690058940 
  H   -4.91574059417615      1.39410577352380      0.56084134748227 
  H   -4.94171318776417     -1.36039101945015      0.41814394547521 
  H   -0.27366218218347     -3.65684137259083      4.10188738919502 
  H   0.31966790105189      3.72258782672313      3.94755548354672 
 
[CuII(NHC2)]2+ 
Charge = 2+, Spin Multiplicity = 2 
  Cu  0.10082553881777     -0.01162902283546      0.01361636182430 
  N   0.09177864474464     -2.19517970873212      0.32308094264054 
  N   0.09137194740309      2.19001728077502      0.10570309438298 
  N   0.10827421165350     -0.07099431073676     -2.35899804452316 
  C   2.01124054044758     -0.01428398467647      0.28199763732638 
  C   -1.82323092576793      0.01339166647361      0.11270973440238 
  C   1.18034605791675     -2.91675405129246     -0.00751777317502 
  C   -0.99026703699657     -2.83143876581840      0.80973520770289 
  C   -0.93174111333120      2.86976818341817     -0.44209041501289 
  C   1.14371215887214      2.87334360283847      0.59769502669000 
  C   -0.11718869170953      0.02223724123049     -3.48195327172126 
  N   2.84037624779476     -1.02856194168101      0.00799105287700 
  N   2.77700097053876      0.95431556545254      0.80012869398815 
  N   -2.61612552028967      0.99388483820824     -0.34147234682676 
  N   -2.64373123466321     -0.91174734175171      0.62498812702716 
  C   2.38269983086823     -2.24831526300419     -0.64617296499871 
  C   1.21761556807379     -4.29737019646537      0.14632964023301 
  C   -1.02253367749143     -4.21586666791314      0.94268663194138 
  C   -2.17909959104891     -2.07492271707419      1.36978218695061 
  C   -2.10415391943216      2.14384708407699     -1.07304378293417 
  C   -0.92461996202856      4.25766057168222     -0.52933508712561 
  C   1.21860276702133      4.25803486249495      0.50658027466162 
  C   2.22929371133160      2.17144489220797      1.38788909987281 
  C   -0.41604792704665      0.15140308518597     -4.89588906257067 
  C   4.14517063482166     -0.70555602282581      0.35413558928361 
  C   4.10388707112358      0.54935114347680      0.86280404025582 
  C   -3.95189824633246      0.68780182495195     -0.12024371896655 
  C   -3.96936544630035     -0.51522370940242      0.50154746212983 
  C   0.09689507503996     -4.96063801549148      0.61736188638752 
  C   0.17302909058302      4.96230233947447     -0.06662095740645 
  H   3.21066160351901     -2.95194856739489     -0.64678144087105 
  H   2.13895157561739     -2.00919453198113     -1.68446866059659 
  H   2.11412694080692     -4.84477824989436     -0.11005708341445 
  H   -1.91335528157941     -4.69984695798052      1.31905774454038 


--- Page 36 ---
 
36 
  H   -3.01551588041271     -2.76487493900858      1.44658839809787 
  H   -1.92902229617680     -1.75297190112432      2.38566071904375 
  H   -2.92268887729186      2.84834785074471     -1.19330006104934 
  H   -1.80378410783041      1.80682969024233     -2.06775743916666 
  H   -1.76878088519060      4.77563038127525     -0.96352176834137 
  H   2.08386636337821      4.77757918354921      0.89433559223726 
  H   3.05339582206838      2.86444222148538      1.53481952833578 
  H   1.83103037585415      1.92599991430796      2.37701918926787 
  H   -0.78462089833580     -0.79806730425287     -5.28657247832320 
  H   0.48137764796439      0.44315832432087     -5.44263359472435 
  H   -1.18177964697345      0.91317245703750     -5.04664450917573 
  H   4.97275893138423     -1.37640146657256      0.20585017751783 
  H   4.88667044492290      1.17136807304030      1.25885312452422 
  H   -4.75715137086282      1.33490899090250     -0.41987965305984 
  H   -4.79145713149439     -1.10545798812091      0.86591334386195 
  H   0.09857593095695     -6.03605126121241      0.73575290932416 
  H   0.21004011506214      6.04089753838973     -0.14160812734517 
 
[CuIII(NHC2)]3+ 
Charge = 3+, Spin Multiplicity = 1 
  Cu  -0.07557324534937      0.08473907599887      0.01749033748120 
  N   -0.11710427142914     -1.92599544467153      0.18499604266301 
  N   -0.03408855310710      2.09834319628169      0.15192748564677 
  N   -0.06374217098734      0.01037443838644     -2.35450306979143 
  C   -1.96331455466345      0.09258358873240      0.11688310543760 
  C   1.81096597131832      0.07711566686513      0.16051460223402 
  N   -2.77232298670605     -0.82451915396435     -0.40362265527385 
  N   -2.72203492841614      0.98952662489733      0.73832089974582 
  N   2.63204382830595      0.97882855247198     -0.36652835933578 
  N   2.55228490983997     -0.79376132700535      0.83765055387342 
  C   0.87877312657885     -2.59424048389109      0.82659624639028 
  C   -1.09389580910824     -2.63335378563464     -0.43741378792053 
  C   -4.09293111422718     -0.50350057690443     -0.11784575943584 
  C   -2.25768058060771     -1.96631986720113     -1.13886379999260 
  C   -4.06096674973948      0.64521354291859      0.60729739372831 
  C   -2.13349050179208      2.06810777191251      1.52145382510592 
  C   -1.02672185965424      2.78103419215742      0.78058873325227 
  C   0.97591682431046      2.79223106044453     -0.43609813244297 
  C   2.11844492599363      2.09608390253589     -1.14279545936202 
  C   3.94375593143155      0.67702411647537     -0.02512379657418 
  C   3.89331171515817     -0.44723814095791      0.73785917581180 
  C   1.93263293339283     -1.85756569651206      1.61918072477104 
  C   -0.09402854849986     -0.10324945645329     -3.49935864804762 


--- Page 37 ---
 
37 
  C   0.93675857382873     -3.97870713513175      0.81292299897897 
  C   -1.05987802315533     -4.02119018052970     -0.47321905287160 
  C   -1.03533297770752      4.16774288692284      0.81203754717493 
  C   0.99323556893107      4.17966249677476     -0.42356411948387 
  C   -0.14275362167436     -0.25197924607384     -4.93928720316056 
  C   -0.02976206203020     -4.70955078177724      0.14198565773654 
  C   -0.02058348487005      4.88439121350562      0.20157150153427 
  H   -4.92310351946544     -1.10006077819544     -0.45629818867217 
  H   -1.95560604924565     -1.63971717101380     -2.13446042588483 
  H   -3.06023219546215     -2.69003245669721     -1.24701954315383 
  H   -4.85678960071782      1.22669557552605      1.04110113890439 
  H   -1.75359426036551      1.65753061436998      2.46168888551861 
  H   -2.91483345337440      2.78007962827257      1.76997881816347 
  H   1.79037051401078      1.72945201992668     -2.11654112956889 
  H   2.92484818011919      2.80617806828481     -1.29933571930475 
  H   4.78209301368428      1.26733943187705     -0.35442103291303 
  H   4.67697464507442     -1.00891539796587      1.21761859523896 
  H   1.49438532830543     -1.42573816220245      2.52335930794152 
  H   2.70413584773577     -2.55494548654925      1.93135083606764 
  H   1.74464354535760     -4.47883796986312      1.32916578102452 
  H   -1.85262374946789     -4.55439654421655     -0.98058415894477 
  H   -1.84120594589389      4.67892136199888      1.32084502853338 
  H   1.81370984554407      4.70073799760100     -0.89806255753123 
  H   -0.88708120023241     -1.00215103290511     -5.21138561163383 
  H   0.83174306785771     -0.57077220076510     -5.31418354922233 
  H   -0.41730287074995      0.69834839761058     -5.40116043977222 
  H   0.01265717110420     -5.79086753088648      0.11307512896297 
  H   -0.01545065818304      5.96677143621972      0.22140529737308 
 
 


--- Page 38 ---
 
38 
[CuIII(NHC4)]3+ 
Charge = 3+, Spin Multiplicity = 1 
  Cu  0.03229808776766     -0.00802812810694      0.11864237027951 
  N   -2.72236441013310     -1.01822590913991     -0.14960253623219 
  N   -2.70775966645164      1.05165990371497      0.36015140049886 
  N   -0.99182677179212      2.73663661758113      0.42799817821532 
  N   1.10074448186391      2.73600010808165      0.01540624132707 
  N   2.78810731813870      1.01965621725733     -0.04001317343073 
  N   2.76758880744168     -1.08680411723597      0.29827934168672 
  N   1.05070082187890     -2.77168505299039      0.22578239216403 
  N   -1.03708489976242     -2.73351942313217     -0.20703421110259 
  C   -1.89907920868679      0.01048829559839      0.10758493656803 
  C   -4.05775312012779     -0.63485251059755     -0.05570934873197 
  C   -4.04869988195395      0.68771926887409      0.26135656080057 
  C   -5.19138504167512     -1.57367244889198     -0.25897594410857 
  H   -6.12600910882298     -1.06898772618353     -0.02356893674340 
  H   -5.10536620157276     -2.43969702346000      0.40068144816933 
  H   -5.24528614831799     -1.93546636579990     -1.28849847318689 
  C   -5.17009958081677      1.64258438237474      0.45740994034379 
  H   -6.11029521854440      1.15027837403222      0.21767312335027 
  H   -5.06873379092374      2.50662916473243     -0.20273830230252 
  H   -5.22470287823333      2.00610692132287      1.48623488531919 
  C   -2.31219154250611      2.34561998129228      0.87111040795094 
  H   -2.33709253692771      2.32833531268671      1.96202245862370 
  H   -3.02204397616674      3.08265205266051      0.51020289364037 
  C   0.04929811497074      1.92107752399187      0.19712738025223 
  C   -0.60614421139231      4.07400698785271      0.38639837453204 
  C   0.72966576443080      4.07370170489001      0.13124650642465 
  C   -1.55511458342522      5.20210040963273      0.57063851288606 
  H   -1.04222487275064      6.14195952138863      0.37751493107682 
  H   -1.96282940532215      5.23351241760329      1.58376877477134 
  H   -2.39064239485964      5.12937486658476     -0.12870112140424 
  C   1.68813126636855      5.20216727639866      0.00619637728472 
  H   1.18240771468055      6.13518675668267      0.24578176616021 
  H   2.09832373591545      5.28245937720486     -1.00309526977729 
  H   2.52151920070277      5.08699086027251      0.70269328150244 
  C   2.41612631650466      2.35588877618411     -0.45226475926732 
  H   2.44440679944015      2.41932726919842     -1.54145435455628 
  H   3.13762346138938      3.05191888289775     -0.03670258513487 
  C   1.96242168578846     -0.02585449938165      0.12790937212463 
  C   4.12236796040711      0.62549381989688      0.02620426476268 
  C   4.10933926116591     -0.71846057305062      0.23343204507798 
  C   5.26123619714263      1.57192141960304     -0.09399483575030 


--- Page 39 ---
 
39 
  H   6.19205370099001      1.04611167414192      0.10763270649858 
  H   5.17415179500235      2.38419019165571      0.63054019690890 
  H   5.32398484517666      2.01323409360016     -1.09136045474970 
  C   5.22958594875976     -1.68700759949008      0.35361922666727 
  H   6.17097695510800     -1.17776251116382      0.15862651447087 
  H   5.13053328307037     -2.49341065413460     -0.37621741315490 
  H   5.27957754154770     -2.13449804127892      1.34893919740984 
  C   2.36828377879116     -2.41610887998863      0.70653077817220 
  H   2.38602835536991     -2.48132685080366      1.79578885642928 
  H   3.08034367616995     -3.12447702349841      0.29596120001931 
  C   0.01598945350023     -1.93740482798923      0.03696398141756 
  C   0.65782363088973     -4.10203986401947      0.09946190899722 
  C   -0.67541172759528     -4.07776919045526     -0.16798302768235 
  C   1.59657171863375     -5.24699527430679      0.22480226222433 
  H   1.07710729473792     -6.17042463433118     -0.02226232696780 
  H   1.99983599748448     -5.33883901388458      1.23589825924429 
  H   2.43561699786388     -5.14185993409134     -0.46630372824764 
  C   -1.64163317753147     -5.18853478484637     -0.36743030588269 
  H   -1.14780280834696     -6.13731469890547     -0.16818145322813 
  H   -2.03474670552326     -5.21195553448894     -1.38662598982001 
  H   -2.48586625839380     -5.10223468886956      0.31952199784533 
  C   -2.34451243479010     -2.31620510404550     -0.66324650498473 
  H   -3.07156508050730     -3.04201412473838     -0.31438641450005 
  H   -2.35692955924031     -2.29174160558918     -1.75433040715131 
 
[CuIII(CF3)4]− 
Charge = 1-, Spin Multiplicity = 1 
  C   -0.57888321333209     -0.87192810770068     -4.61669655249915 
  C   1.24870669476421     -3.05020976763769     -4.40963955185390 
  C   1.47833807231003      0.87018718488180     -3.72877744788674 
  C   3.37418382159552     -1.11507810433457     -4.53219278742627 
  Cu  1.38496409074348     -1.05037661636376     -4.31751573386117 
  F   3.87088529843570      0.02400621230627     -5.10562741861118 
  F   3.79458304414789     -2.10584524235780     -5.37497292534176 
  F   4.08208039740395     -1.28693565968239     -3.38227260075224 
  F   2.51557920676023      1.11013329617410     -2.87255884171744 
  F   1.61325680072138      1.77897519107215     -4.73225108258272 
  F   0.37859040641269      1.26585090019746     -3.01888290997497 
  F   -1.08831195537215     -1.80859502505455     -5.47048544593969 
  F   -0.89399280688326      0.31876306969050     -5.21409482456034 
  F   -1.34336207842653     -0.94156858200232     -3.49436571889668 
  F   0.13488502680767     -3.54372229366747     -3.78682883279539 
  F   1.22150934303065     -3.56228088779489     -5.66753459725558 


--- Page 40 ---
 
40 
  F   2.27957453288055     -3.68768122272611     -3.77795042304449 
\newpage

\end{document}